%% Section 1, Introduction. Converted from paper/sections/01_introduction.md.
\section{Introduction}
\label{sec:intro}

A hand is dexterous when it can change an object's pose without putting the object down. That
property, and not the finger count, is what separates a hand from a gripper. Bicchi
\cite{bicchi_grasping_chapter_2001} draws the same line in Sec.~1.1, between restraining an object
and ``manipulating objects with fingers, in contrast to manipulation with the robot arm''. Restraint
is a static question about whether the contacts prevent motion. In-hand manipulation is a dynamic
question about whether contacts can be broken and remade while the object stays held. An et al.
\cite{an_dexil_survey_2025} put the same idea in their abstract as the ability ``to skillfully
control, reorient, and manipulate objects through precise, coordinated finger movements and adaptive
force modulation''. Both definitions place the work in the fingers rather than in the arm.

The analytic theory answered the static question and stalled on the dynamic one. Form closure has
a first-order test on the grasp matrix and known contact counts, four in the plane and seven in
three dimensions for any polyhedron, per \cite{bicchi_grasping_chapter_2001} Sec.~1.3.1. Force closure
adds the wrench balance and the hand Jacobian. Neither delivers what a controller needs. The same
chapter states in Sec.~1.5 that ``force closure does not guarantee stability'', and its Sec.~1.8
names the reason the theory could not be pushed further, which is that ``the nonsmooth nature of
grasp dynamics, because of the unilateral constraints on displacements and forces, has made a
thorough analysis very difficult''. That verdict is one architect's, on one chapter, and this
survey did not survey the tradition it judges. Learned control did not solve those modelling
problems. It went around them by sampling a simulator instead of solving a model, and by scoring a
rollout instead of certifying a configuration. Of the 112 method papers in this corpus, 53 train
with reinforcement learning and 36 run in Isaac Gym, against 6 on its successor Isaac Lab.

Three things this survey measured are worth stating before the reader commits to 27,000 words. The
first is that papers disagree with their own released code. Sixty-two of the 112 method rows
released code that could be read against the paper, 38 of those record a discrepancy, and nine are
contradictions where the shipped code states a different objective from the published one. The
sharpest case is \cite{physhoi_2023}. Its \texttt{compute\_humanoid\_reward} hardcodes the object rotation and
rotation-velocity errors to zero, with the real computation commented out beside them, while the
reward table in its own paper lists weights of 0.1 and 0.01 for exactly those terms on GRAB. Its
position-only success criterion could not have caught that, and its headline 95.4 percent is cited
as a baseline. Section 5.8 classifies all 38, and seven further accusations an earlier draft made
were withdrawn under adversarial review and recorded beside the charge.
\figurename~\ref{fig:codegap} draws the narrowing.

% ---------------------------------------------------------------------------
% The first finding as a picture: the funnel from 112 method rows to the ten
% whose code contradicts the paper, with the 28 charges the survey withdrew
% drawn at the same weight as the 10 it kept. Written by
% tools/make_finding_figures.py from corpus/rows; the guard keeps the document
% compiling until it is there.
% ---------------------------------------------------------------------------
\makeatletter
\@ifundefined{fnLoaded}{\IfFileExists{figs/finding_numbers.tex}{% generated by tools/make_finding_figures.py; do not edit
% Every count the finding figures and their captions print, recomputed from corpus/rows.
\providecommand{\fnLoaded}{}
\providecommand{\fnRows}{112}
\providecommand{\fnCode}{62}
\providecommand{\fnDisagree}{38}
\providecommand{\fnKept}{9}
\providecommand{\fnOtherClass}{29}
\providecommand{\fnParse}{13}
\providecommand{\fnAbsent}{8}
\providecommand{\fnSkew}{4}
\providecommand{\fnInternal}{4}
\providecommand{\fnPenSettled}{96}
\providecommand{\fnPenHandled}{11}
\providecommand{\fnPenClosed}{4}
\providecommand{\fnPenOther}{7}
\providecommand{\fnPenRollout}{0}
\providecommand{\fnTrials}{39}
\providecommand{\fnTrialsMedian}{15}
\providecommand{\fnTrialsMode}{10}
\providecommand{\fnTrialsModeN}{16}
\providecommand{\fnTotals}{24}
\providecommand{\fnTotalsMin}{12}
\providecommand{\fnTotalsMax}{750}
\providecommand{\fnUnclear}{7}
}{}}{}
\makeatother
\begin{figure}[!t]
  \centering
  \IfFileExists{figs/fig_codegap.tex}{\resizebox{\columnwidth}{!}{% generated by tools/make_finding_figures.py; do not edit
\begin{tikzpicture}[
  font=\footnotesize,
  bar/.style={draw=none,fill=black!62},
  barmid/.style={draw=none,fill=black!42},
  barlight/.style={draw=none,fill=black!20},
  baracc/.style={draw=none,fill=accent},
  seg/.style={draw=white,line width=0.5pt},
  lnk/.style={draw=black!38,line width=0.28pt},
  ghost/.style={draw=black!45,line width=0.3pt,dash pattern=on 1.2pt off 1.2pt},
  lbl/.style={font=\scriptsize},
  small/.style={font=\scriptsize\color{black!55}},
  num/.style={font=\scriptsize\color{black!55}},
]

\fill[barlight] (0.00,0.00) rectangle (4.90,-0.52);
\node[font=\small\bfseries\color{black!75}] at (2.45,-0.26) {112};
\node[lbl,anchor=west,align=left] at (5.20,-0.26) {method papers\\in the corpus};
\fill[black!7] (0.00,-0.52) -- (4.90,-0.52) -- (3.81,-0.78) -- (1.09,-0.78) -- cycle;
\fill[barmid] (1.09,-0.78) rectangle (3.81,-1.30);
\node[font=\small\bfseries\color{white}] at (2.45,-1.04) {62};
\node[lbl,anchor=west,align=left] at (5.20,-1.04) {released code that could\\be read against the paper};
\fill[black!7] (1.09,-1.30) -- (3.81,-1.30) -- (3.28,-1.56) -- (1.62,-1.56) -- cycle;
\fill[bar] (1.62,-1.56) rectangle (3.28,-2.08);
\node[font=\small\bfseries\color{white}] at (2.45,-1.82) {38};
\node[lbl,anchor=west,align=left] at (5.20,-1.82) {disagree with their\\own code somehow};
\fill[black!5] (1.62,-2.08) -- (3.28,-2.08) -- (7.60,-2.74) -- (0,-2.74) -- cycle;
\node[anchor=south east,font=\scriptsize\color{black!45}] at (7.60,-2.68) {one mark, one paper};
\fill[baracc] (0.00,-2.74) rectangle (0.15,-3.16);
\fill[baracc] (0.20,-2.74) rectangle (0.35,-3.16);
\fill[baracc] (0.40,-2.74) rectangle (0.55,-3.16);
\fill[baracc] (0.60,-2.74) rectangle (0.75,-3.16);
\fill[baracc] (0.80,-2.74) rectangle (0.95,-3.16);
\fill[baracc] (1.00,-2.74) rectangle (1.15,-3.16);
\fill[baracc] (1.20,-2.74) rectangle (1.35,-3.16);
\fill[baracc] (1.40,-2.74) rectangle (1.55,-3.16);
\fill[baracc] (1.60,-2.74) rectangle (1.75,-3.16);
\fill[barlight] (1.80,-2.74) rectangle (1.95,-3.16);
\fill[barlight] (2.00,-2.74) rectangle (2.15,-3.16);
\fill[barlight] (2.20,-2.74) rectangle (2.35,-3.16);
\fill[barlight] (2.40,-2.74) rectangle (2.55,-3.16);
\fill[barlight] (2.60,-2.74) rectangle (2.75,-3.16);
\fill[barlight] (2.80,-2.74) rectangle (2.95,-3.16);
\fill[barlight] (3.00,-2.74) rectangle (3.15,-3.16);
\fill[barlight] (3.20,-2.74) rectangle (3.35,-3.16);
\fill[barlight] (3.40,-2.74) rectangle (3.55,-3.16);
\fill[barlight] (3.60,-2.74) rectangle (3.75,-3.16);
\fill[barlight] (3.80,-2.74) rectangle (3.95,-3.16);
\fill[barlight] (4.00,-2.74) rectangle (4.15,-3.16);
\fill[barlight] (4.20,-2.74) rectangle (4.35,-3.16);
\fill[barlight] (4.40,-2.74) rectangle (4.55,-3.16);
\fill[barlight] (4.60,-2.74) rectangle (4.75,-3.16);
\fill[barlight] (4.80,-2.74) rectangle (4.95,-3.16);
\fill[barlight] (5.00,-2.74) rectangle (5.15,-3.16);
\fill[barlight] (5.20,-2.74) rectangle (5.35,-3.16);
\fill[barlight] (5.40,-2.74) rectangle (5.55,-3.16);
\fill[barlight] (5.60,-2.74) rectangle (5.75,-3.16);
\fill[barlight] (5.80,-2.74) rectangle (5.95,-3.16);
\fill[barlight] (6.00,-2.74) rectangle (6.15,-3.16);
\fill[barlight] (6.20,-2.74) rectangle (6.35,-3.16);
\fill[barlight] (6.40,-2.74) rectangle (6.55,-3.16);
\fill[barlight] (6.60,-2.74) rectangle (6.75,-3.16);
\fill[barlight] (6.80,-2.74) rectangle (6.95,-3.16);
\fill[barlight] (7.00,-2.74) rectangle (7.15,-3.16);
\fill[barlight] (7.20,-2.74) rectangle (7.35,-3.16);
\fill[barlight] (7.40,-2.74) rectangle (7.55,-3.16);
\draw[lnk] (0,-3.22) -- (7.60,-3.22);
\node[anchor=north west,align=left,font=\scriptsize\color{accent}] at (0,-3.30) {\textbf{9} contradictions: the shipped code states a different objective};
\node[anchor=north west,align=left,font=\scriptsize\color{black!58}] at (0,-3.56) {\textbf{29} of another kind, classified and not retracted: 13 a limit of this\\survey's own parse, 8 a component never released, 4 version skew, 4 a paper\\disagreeing with itself};
\end{tikzpicture}
}}{%
    \framebox[0.95\columnwidth]{\rule{0pt}{2.2cm}\textit{figs/fig\_codegap.tex pending}}}
  \caption{Papers disagree with their own released code, and most of the disagreements are this
  survey's own limits rather than theirs. Of the \fnRows{} method rows, \fnCode{} released code that
  could be read against the paper and \fnDisagree{} of those record a disagreement. The lower bar
  reopens those \fnDisagree{} at their own scale: \fnKept{} are contradictions, where the shipped
  code states a different objective from the published one, and \fnWithdrawn{} are withdrawn, most
  of them about this survey's own parse or about what a repository shipped rather than about the
  work. Further accusations
  an earlier draft made were withdrawn under adversarial review, each recorded in the accused row's
  \texttt{mismatch\_review} field beside the charge. Section~\ref{sec:train:code} classifies all
  \fnDisagree{} and names them.}
  \label{fig:codegap}
\end{figure}

The second is that the quantity most specific to a hand is the one closed-loop policies do not
record. Eleven of the 96 method rows whose notes settle the question address interpenetration at all,
seven of the eleven do it outside a closed-loop policy in a grasp synthesiser, a trajectory optimiser
or a contact model, and we found none that reports a penetration number for its own trained policy's
rollouts. The claim is about learned closed-loop control and not about the field. Grasp synthesis and
hand-object reconstruction have reported penetration depth and intersection volume as comparative
columns for years, and four rows of this corpus do it: OakInk \cite{oakink_2022} tabulates
penetration depth, solid intersection volume and simulation displacement across methods, BiDexGrasp
\cite{bidexgrasp_2026} prints penetration depth beside a prior method's, BimanGrasp
\cite{bimangrasp_2024} fails any grasp whose total penetration exceeds 1.5 mm, and TopoRetarget
\cite{toporetarget_2026} reports a maximum penetration and a share of frames past 2 mm against a
baseline retargeter. Every one of those numbers scores a pose or a reference trajectory rather than
the behaviour a trained policy produced, and it is the rollout that is missing. The
obstacle is not the engines. NVIDIA's own IsaacGymEnvs repository already computes a
per-environment maximum interpenetration depth in Warp and gates the policy update on a 1 mm
threshold. Section 7.3 has the file and the lines.

The third is that hardware and published work have come apart. 18 of the 33 hand rows in
Table~\ref{tab:hands-available} and Table~\ref{tab:hands-announced}
appear in no method row, and 7 of those can be bought today or built from
published designs. Thirty-five method rows run on the Allegro, whose weight, joint torque, payload
and price have no reachable source, because its product page returns HTTP 404 and everything
Table~\ref{tab:hands-available} confirms about it comes from its ROS driver.

None of those three findings is a first, and the audit behind the first of them is not a new idea.
Collberg and Proebsting \cite{collberg_repeatability_2016} examined 601 papers in computer systems
research for whether the code behind them could be obtained and built at all. Xu et al.
\cite{biocon_2026} align 48 bioinformatics projects with their publications at sentence-to-function
granularity under expert annotation, and SciCoQA \cite{scicoqa_2026} collects 92 real paper-code
discrepancies, mined from issue trackers and reproducibility reports, into a benchmark for detecting
such discrepancies automatically. In reinforcement learning the phenomenon itself is a known result.
Engstrom et al. \cite{engstrom_implementation_matters_2020} show that code-level optimisations
present only in the implementation account for most of PPO's reported gain over TRPO, and Meta-World+
\cite{metaworld_plus_2025} finds undocumented changes accumulated across one benchmark's own
versions, which make comparisons between those versions unfair. Both establish it on a single
codebase. The closest relative to the audit here is Knox et al. \cite{knox_reward_misdesign_2023},
who review nineteen reinforcement-learning publications on autonomous driving, characterise the
reward functions of ten of them exhaustively in a standard form, apply eight sanity checks and report
near-universal flaws in reward design. Their ground truth for what each reward was is the authors,
obtained through correspondence with them rather than by reading a released repository, and what they
establish is that published reward descriptions are incomplete: one of the ten described its reward,
discount factor, termination conditions and timestep thoroughly. Reading the code instead needs no
correspondence and supports a different charge, which is that where code exists it sometimes
contradicts the description. Raff \cite{raff_reproducibility_2019} took the opposite ground truth on
purpose, reimplementing 255 papers from their text alone and never opening the authors' code, which
is what makes the choice of arbiter a position rather than an accident.

So the claim here is narrow. We are aware of no prior work in robotics, and none in dexterous
manipulation, that reads a field's released reward implementations against the rewards its own papers
describe. Those seven comparisons are cited from outside this corpus and enter none of its counts.
The exposure of the claim belongs beside the count above: the evidence is a repository at a fetched
commit, each hash recorded in \texttt{corpus/code\_manifest.json}, and a repository at a commit is
evidence about that repository rather than about the run that produced a paper's numbers, because the
commit may postdate, precede or diverge from it. Hora \cite{hora_2022} states the problem in its own
README, which directs a reader to tag v0.0.1 and not to the commit parsed here to reproduce the
published numbers. The seven accusations an earlier draft withdrew are recorded in the accused rows
for the same reason: they are the evidence that the charges left standing were checked rather than
counted.

Fourteen corpus entries are themselves surveys or engine-comparison studies. Appendix D says what
each covers, on one set of columns tabulated in the repository. Four of the fourteen could
not be obtained, or were fetched too late to read, and are entered as such. Two of the three things
this survey adds are columns none of the fourteen fills. The first is
Table~\ref{tab:simulators}, which takes the engines \cite{nine_physics_engines_review_2024} scored
on documentation and usability and adds contact model, solver, iteration count, default timestep and
penetration exposure, conditioned on what a hand does to a solver. Erez et al.
\cite{physics_engine_comparison_2015} and Le Lidec et al. \cite{contact_models_comparison_2023} do
measure engines, on five engines and on four contact formulations, and neither surveys the field
those engines are used in. The second is two hands on one object as its own problem, which none of
the four field surveys gives more than a subsection.

The third is the penetration measurement, on an axis a predecessor had already named. Zhao et al.
\cite{zhao_dexhand_survey_2026} state in their Sec.~IV-C that the field assesses ``at least two
layers of performance: the quality of grasps or poses prior to execution, and the performance of
policies or generators during downstream execution'', and names ``physical plausibility, including
penetration'' as the most common criterion for the first layer. That is the reference-versus-rollout
split, named before this survey measured it, and the idea that contact quality belongs on the
evaluation axis is Zhao's. What Sec.~IV-C does not give is a measurement method or a count, and this
survey supplies those two. It does not supply a threshold. The 2 mm figure the field uses is taken
from \cite{toporetarget_2026} with no independent justification, the captured human grasps in
\cite{grab_2020} sit above it at 3.25 mm, and Section 7.7 states plainly that 2 mm is a simulator
convention rather than a physical bound.

The scope is single-hand and bimanual multi-fingered manipulation. Parallel-jaw manipulation
enters only as a comparison, which matters because 12 of the 112 method rows name a parallel-jaw
gripper among their own embodiments. Locomotion is excluded. Prosthetics are excluded except where
a hand crosses over into robot use, as the Psyonic Ability Hand does. The corpus holds 221
bibliography entries, of which 218 carry a structured row read from a note. Those 218 rows are 112
method papers, 33 hands, 15 simulators, 15 datasets, 14 benchmarks, 14 surveys, 8 tactile sensors
and 7 evaluation protocols. ``Method row'' throughout means one of the 112, and every headline count
in this survey has one of those eight classes as its denominator.

Two limits apply to every number here. This is a corpus of the learned era, which is a selection
effect and not a judgement: 138 of the 221 entries are dated 2024 or later, 16 predate 2018, and
two of the 112 method rows predate 2018, so any statement about a trend over time is a statement
about 2022 onward. The analytic tradition is represented by one readable chapter rather than
surveyed, and the planning line that took up the dynamic question directly, finger gaiting and
rolling-contact manipulation and regrasp planning, is not in the corpus at all. Every coverage
statistic is also a floor rather than a rate, because a value this survey failed to extract is
indistinguishable from a value the paper never reported, and every such miss converts a reporting
paper into a silent one.

\figurename~\ref{fig:field} puts the field on one page. Section 2 sets out the task families and what makes each
hard. Section 3 covers hands, their makers, and the gap between what is sold and what is run.
Section 4 covers simulators and contact models. Section 5 covers how policies are trained. Section
6 covers bimanual work as its own problem. Section 7 proposes an evaluation frame rather than a
leaderboard, and is the longest section. Section 8 states the gaps as claims with their evidence.
Section 9 says what to do about them, addressed to someone publishing, running experiments or
buying a hand. Appendix A is the method, and says where the tabulation this survey is built on
lives, which is the repository rather than these pages. Appendix B carries the three tables
Sections 3 and 4 read cells out of, Appendix C is the reward extraction, and Appendix D is the
comparison with the existing surveys.

% ---------------------------------------------------------------------------
% Figure 1, the field on one page. Full width: five columns of nodes with the
% routes that exist in the corpus. Drawn by tools/make_tikz_figures.py into
% figs/fig_field.tex; the guard keeps the document compiling until it is there.
% ---------------------------------------------------------------------------
\begin{figure*}[!t]
  \centering
  \IfFileExists{figs/fig_field.tex}{% generated by tools/make_tikz_figures.py; do not edit
\begin{tikzpicture}[
  font=\scriptsize,
  bar/.style={draw=none,fill=black!62},
  barlight/.style={draw=none,fill=black!22},
  baracc/.style={draw=none,fill=accent},
  box/.style={draw=black!45,line width=0.3pt,inner sep=3pt,align=left},
  ghost/.style={draw=black!35,line width=0.3pt,dash pattern=on 1.4pt off 1.4pt,inner sep=3pt,align=left},
  lnk/.style={draw=black!38,line width=0.28pt},
  lbl/.style={font=\scriptsize},
  num/.style={font=\scriptsize\color{black!55}},
]

\draw[draw=black!41,line width=0.44pt] (2.64,-3.15) .. controls (3.19,-3.15) and (3.29,-6.15) .. (3.84,-6.15);
\draw[draw=black!41,line width=0.44pt] (6.48,-6.38) .. controls (7.03,-6.38) and (7.13,-5.70) .. (7.68,-5.70);
\draw[draw=black!41,line width=0.44pt] (10.32,-3.32) .. controls (10.87,-3.32) and (10.97,-6.42) .. (11.52,-6.42);
\draw[draw=black!41,line width=0.44pt] (10.32,-5.57) .. controls (10.87,-5.57) and (10.97,-5.73) .. (11.52,-5.73);
\draw[draw=black!41,line width=0.44pt] (10.32,-5.73) .. controls (10.87,-5.73) and (10.97,-7.88) .. (11.52,-7.88);
\draw[draw=black!41,line width=0.44pt] (10.32,-6.32) .. controls (10.87,-6.32) and (10.97,-3.14) .. (11.52,-3.14);
\draw[draw=black!41,line width=0.44pt] (10.32,-6.59) .. controls (10.87,-6.59) and (10.97,-6.71) .. (11.52,-6.71);
\draw[draw=black!41,line width=0.44pt] (10.32,-7.04) .. controls (10.87,-7.04) and (10.97,-1.78) .. (11.52,-1.78);
\draw[draw=black!41,line width=0.44pt] (10.32,-7.20) .. controls (10.87,-7.20) and (10.97,-3.31) .. (11.52,-3.31);
\draw[draw=black!41,line width=0.44pt] (14.16,-2.77) .. controls (14.71,-2.77) and (14.81,-3.56) .. (15.36,-3.56);
\draw[draw=black!41,line width=0.44pt] (14.16,-3.02) .. controls (14.71,-3.02) and (14.81,-4.68) .. (15.36,-4.68);
\draw[draw=black!41,line width=0.44pt] (14.16,-6.57) .. controls (14.71,-6.57) and (14.81,-4.93) .. (15.36,-4.93);
\draw[draw=black!41,line width=0.44pt] (14.16,-7.52) .. controls (14.71,-7.52) and (14.81,-4.02) .. (15.36,-4.02);
\draw[draw=black!42,line width=0.48pt] (2.64,-5.98) .. controls (3.19,-5.98) and (3.29,-3.22) .. (3.84,-3.22);
\draw[draw=black!42,line width=0.48pt] (2.64,-6.38) .. controls (3.19,-6.38) and (3.29,-6.38) .. (3.84,-6.38);
\draw[draw=black!42,line width=0.48pt] (6.48,-2.81) .. controls (7.03,-2.81) and (7.13,-5.29) .. (7.68,-5.29);
\draw[draw=black!42,line width=0.48pt] (6.48,-5.46) .. controls (7.03,-5.46) and (7.13,-6.53) .. (7.68,-6.53);
\draw[draw=black!42,line width=0.48pt] (6.48,-6.15) .. controls (7.03,-6.15) and (7.13,-1.85) .. (7.68,-1.85);
\draw[draw=black!42,line width=0.48pt] (10.32,-4.53) .. controls (10.87,-4.53) and (10.97,-5.47) .. (11.52,-5.47);
\draw[draw=black!42,line width=0.48pt] (10.32,-4.65) .. controls (10.87,-4.65) and (10.97,-6.57) .. (11.52,-6.57);
\draw[draw=black!42,line width=0.48pt] (10.32,-5.42) .. controls (10.87,-5.42) and (10.97,-2.98) .. (11.52,-2.98);
\draw[draw=black!42,line width=0.48pt] (10.32,-6.46) .. controls (10.87,-6.46) and (10.97,-4.58) .. (11.52,-4.58);
\draw[draw=black!42,line width=0.48pt] (10.32,-7.36) .. controls (10.87,-7.36) and (10.97,-6.85) .. (11.52,-6.85);
\draw[draw=black!42,line width=0.48pt] (14.16,-4.35) .. controls (14.71,-4.35) and (14.81,-4.80) .. (15.36,-4.80);
\draw[draw=black!43,line width=0.52pt] (10.32,-1.35) .. controls (10.87,-1.35) and (10.97,-3.90) .. (11.52,-3.90);
\draw[draw=black!43,line width=0.52pt] (10.32,-2.85) .. controls (10.87,-2.85) and (10.97,-2.64) .. (11.52,-2.64);
\draw[draw=black!43,line width=0.52pt] (10.32,-4.40) .. controls (10.87,-4.40) and (10.97,-4.35) .. (11.52,-4.35);
\draw[draw=black!43,line width=0.52pt] (10.32,-4.78) .. controls (10.87,-4.78) and (10.97,-7.70) .. (11.52,-7.70);
\draw[draw=black!43,line width=0.52pt] (10.32,-5.26) .. controls (10.87,-5.26) and (10.97,-1.26) .. (11.52,-1.26);
\draw[draw=black!43,line width=0.52pt] (14.16,-7.70) .. controls (14.71,-7.70) and (14.81,-5.06) .. (15.36,-5.06);
\draw[draw=black!44,line width=0.56pt] (6.48,-3.35) .. controls (7.03,-3.35) and (7.13,-7.20) .. (7.68,-7.20);
\draw[draw=black!44,line width=0.56pt] (10.32,-3.55) .. controls (10.87,-3.55) and (10.97,-7.52) .. (11.52,-7.52);
\draw[draw=black!44,line width=0.56pt] (10.32,-4.27) .. controls (10.87,-4.27) and (10.97,-2.81) .. (11.52,-2.81);
\draw[draw=black!45,line width=0.60pt] (2.64,-6.18) .. controls (3.19,-6.18) and (3.29,-5.35) .. (3.84,-5.35);
\draw[draw=black!45,line width=0.60pt] (6.48,-3.08) .. controls (7.03,-3.08) and (7.13,-6.25) .. (7.68,-6.25);
\draw[draw=black!45,line width=0.60pt] (10.32,-1.95) .. controls (10.87,-1.95) and (10.97,-7.35) .. (11.52,-7.35);
\draw[draw=black!45,line width=0.60pt] (10.32,-4.14) .. controls (10.87,-4.14) and (10.97,-1.00) .. (11.52,-1.00);
\draw[draw=black!45,line width=0.60pt] (14.16,-1.33) .. controls (14.71,-1.33) and (14.81,-4.55) .. (15.36,-4.55);
\draw[draw=black!46,line width=0.64pt] (10.32,-0.95) .. controls (10.87,-0.95) and (10.97,-0.48) .. (11.52,-0.48);
\draw[draw=black!46,line width=0.64pt] (10.32,-1.75) .. controls (10.87,-1.75) and (10.97,-6.28) .. (11.52,-6.28);
\draw[draw=black!46,line width=0.64pt] (10.32,-6.18) .. controls (10.87,-6.18) and (10.97,-1.52) .. (11.52,-1.52);
\draw[draw=black!46,line width=0.64pt] (14.16,-4.13) .. controls (14.71,-4.13) and (14.81,-3.79) .. (15.36,-3.79);
\draw[draw=black!47,line width=0.68pt] (6.48,-4.83) .. controls (7.03,-4.83) and (7.13,-4.65) .. (7.68,-4.65);
\draw[draw=black!47,line width=0.68pt] (6.48,-5.15) .. controls (7.03,-5.15) and (7.13,-5.49) .. (7.68,-5.49);
\draw[draw=black!48,line width=0.72pt] (14.16,-7.88) .. controls (14.71,-7.88) and (14.81,-7.73) .. (15.36,-7.73);
\draw[draw=black!50,line width=0.80pt] (6.48,-2.54) .. controls (7.03,-2.54) and (7.13,-4.27) .. (7.68,-4.27);
\draw[draw=black!51,line width=0.85pt] (14.16,-6.80) .. controls (14.71,-6.80) and (14.81,-7.35) .. (15.36,-7.35);
\draw[draw=black!53,line width=0.93pt] (6.48,-1.99) .. controls (7.03,-1.99) and (7.13,-1.05) .. (7.68,-1.05);
\draw[draw=black!53,line width=0.93pt] (6.48,-4.52) .. controls (7.03,-4.52) and (7.13,-3.38) .. (7.68,-3.38);
\draw[draw=black!53,line width=0.93pt] (10.32,-1.15) .. controls (10.87,-1.15) and (10.97,-2.48) .. (11.52,-2.48);
\draw[draw=black!53,line width=0.93pt] (14.16,-7.35) .. controls (14.71,-7.35) and (14.81,-2.46) .. (15.36,-2.46);
\draw[draw=black!55,line width=0.97pt] (14.16,-6.33) .. controls (14.71,-6.33) and (14.81,-2.07) .. (15.36,-2.07);
\draw[draw=black!56,line width=1.01pt] (10.32,-1.55) .. controls (10.87,-1.55) and (10.97,-5.20) .. (11.52,-5.20);
\draw[draw=black!56,line width=1.01pt] (10.32,-3.09) .. controls (10.87,-3.09) and (10.97,-4.13) .. (11.52,-4.13);
\draw[draw=black!56,line width=1.01pt] (14.16,-0.94) .. controls (14.71,-0.94) and (14.81,-3.34) .. (15.36,-3.34);
\draw[draw=black!57,line width=1.05pt] (14.16,-5.66) .. controls (14.71,-5.66) and (14.81,-6.98) .. (15.36,-6.98);
\draw[draw=black!58,line width=1.09pt] (14.16,-3.90) .. controls (14.71,-3.90) and (14.81,-1.28) .. (15.36,-1.28);
\draw[draw=black!59,line width=1.13pt] (2.64,-5.08) .. controls (3.19,-5.08) and (3.29,-4.83) .. (3.84,-4.83);
\draw[draw=black!59,line width=1.13pt] (14.16,-5.27) .. controls (14.71,-5.27) and (14.81,-1.67) .. (15.36,-1.67);
\draw[draw=black!60,line width=1.17pt] (14.16,-4.58) .. controls (14.71,-4.58) and (14.81,-6.60) .. (15.36,-6.60);
\draw[draw=black!62,line width=1.25pt] (6.48,-4.21) .. controls (7.03,-4.21) and (7.13,-1.45) .. (7.68,-1.45);
\draw[draw=black!62,line width=1.25pt] (14.16,-3.27) .. controls (14.71,-3.27) and (14.81,-6.22) .. (15.36,-6.22);
\draw[draw=black!63,line width=1.29pt] (2.64,-2.11) .. controls (3.19,-2.11) and (3.29,-2.13) .. (3.84,-2.13);
\draw[draw=black!63,line width=1.29pt] (6.48,-2.26) .. controls (7.03,-2.26) and (7.13,-2.79) .. (7.68,-2.79);
\draw[draw=black!67,line width=1.45pt] (14.16,-2.52) .. controls (14.71,-2.52) and (14.81,-0.89) .. (15.36,-0.89);
\draw[draw=black!72,line width=1.61pt] (2.64,-2.63) .. controls (3.19,-2.63) and (3.29,-4.31) .. (3.84,-4.31);
\draw[draw=black!74,line width=1.70pt] (2.64,-4.33) .. controls (3.19,-4.33) and (3.29,-2.67) .. (3.84,-2.67);
\draw[draw=black!74,line width=1.70pt] (10.32,-2.62) .. controls (10.87,-2.62) and (10.97,-0.74) .. (11.52,-0.74);
\draw[draw=black!80,line width=1.94pt] (14.16,-0.54) .. controls (14.71,-0.54) and (14.81,-0.49) .. (15.36,-0.49);
\draw[draw=black!85,line width=2.10pt] (14.16,-1.72) .. controls (14.71,-1.72) and (14.81,-5.84) .. (15.36,-5.84);
\node[font=\footnotesize\bfseries,anchor=south] at (1.32,0.48) {data source};
\node[font=\scriptsize\color{black!60},anchor=south] at (1.32,0.10) {Sec.~\ref{sec:train:human}};
\draw[draw=black!45,line width=0.4pt] (0.00,0.00) -- (2.64,0.00);
\draw[box,fill=white] (0.00,-1.65) rectangle (2.64,-3.61);
\node[lbl,anchor=west,align=left] at (0.10,-2.63) {human data\\named};
\node[num,anchor=east] at (2.54,-2.63) {53};
\draw[box,fill=white] (0.00,-3.77) rectangle (2.64,-5.65);
\node[lbl,anchor=west,align=left] at (0.10,-4.71) {no\\demonstrations};
\node[num,anchor=east] at (2.54,-4.71) {50};
\draw[ghost,fill=white] (0.00,-5.81) rectangle (2.64,-6.55);
\node[lbl,anchor=west,align=left] at (0.10,-6.18) {not stated};
\node[num,anchor=east] at (2.54,-6.18) {9};
\node[font=\footnotesize\bfseries,anchor=south] at (5.16,0.48) {embodiment};
\node[font=\scriptsize\color{black!60},anchor=south] at (5.16,0.10) {Sec.~\ref{sec:hands},~\ref{sec:bimanual}};
\draw[draw=black!45,line width=0.4pt] (3.84,0.00) -- (6.48,0.00);
\draw[box,fill=white] (3.84,-1.65) rectangle (6.48,-3.69);
\node[lbl,anchor=west,align=left] at (3.94,-2.67) {one hand};
\node[num,anchor=east] at (6.38,-2.67) {56};
\draw[box,fill=white] (3.84,-3.85) rectangle (6.48,-5.81);
\node[lbl,anchor=west,align=left] at (3.94,-4.83) {two hands};
\node[num,anchor=east] at (6.38,-4.83) {53};
\draw[ghost,fill=white] (3.84,-5.97) rectangle (6.48,-6.55);
\node[lbl,anchor=west,align=left] at (3.94,-6.26) {not stated};
\node[num,anchor=east] at (6.38,-6.26) {3};
\node[font=\footnotesize\bfseries,anchor=south] at (9.00,0.48) {simulator};
\node[font=\scriptsize\color{black!60},anchor=south] at (9.00,0.10) {Sec.~\ref{sec:simulators}};
\draw[draw=black!45,line width=0.4pt] (7.68,0.00) -- (10.32,0.00);
\draw[ghost,fill=white] (7.68,-0.70) rectangle (10.32,-2.19);
\node[lbl,anchor=west,align=left] at (7.78,-1.45) {not stated};
\node[num,anchor=east] at (10.22,-1.45) {36};
\draw[box,fill=white] (7.68,-2.35) rectangle (10.32,-3.82);
\node[lbl,anchor=west,align=left] at (7.78,-3.09) {Isaac Gym};
\node[num,anchor=east] at (10.22,-3.09) {35};
\draw[box,fill=white] (7.68,-3.98) rectangle (10.32,-4.95);
\node[lbl,anchor=west,align=left] at (7.78,-4.46) {MuJoCo};
\node[num,anchor=east] at (10.22,-4.46) {17};
\draw[box,fill=white] (7.68,-5.11) rectangle (10.32,-5.88);
\node[lbl,anchor=west,align=left] at (7.78,-5.49) {other};
\node[num,anchor=east] at (10.22,-5.49) {10};
\draw[box,fill=white] (7.68,-6.04) rectangle (10.32,-6.73);
\node[lbl,anchor=west,align=left] at (7.78,-6.39) {Isaac Lab or\\Sim};
\node[num,anchor=east] at (10.22,-6.39) {7};
\draw[box,fill=white] (7.68,-6.89) rectangle (10.32,-7.50);
\node[lbl,anchor=west,align=left] at (7.78,-7.20) {SAPIEN};
\node[num,anchor=east] at (10.22,-7.20) {4};
\node[font=\footnotesize\bfseries,anchor=south] at (12.84,0.48) {training paradigm};
\node[font=\scriptsize\color{black!60},anchor=south] at (12.84,0.10) {Sec.~\ref{sec:training}};
\draw[draw=black!45,line width=0.4pt] (11.52,0.00) -- (14.16,0.00);
\draw[box,fill=white] (11.52,-0.15) rectangle (14.16,-2.11);
\node[lbl,anchor=west,align=left] at (11.62,-1.13) {RL};
\node[num,anchor=east] at (14.06,-1.13) {53};
\draw[box,fill=white] (11.52,-2.27) rectangle (14.16,-3.51);
\node[lbl,anchor=west,align=left] at (11.62,-2.89) {BC};
\node[num,anchor=east] at (14.06,-2.89) {27};
\draw[box,fill=white] (11.52,-3.67) rectangle (14.16,-4.81);
\node[lbl,anchor=west,align=left] at (11.62,-4.24) {distillation};
\node[num,anchor=east] at (14.06,-4.24) {23};
\draw[box,fill=white] (11.52,-4.97) rectangle (14.16,-5.96);
\node[lbl,anchor=west,align=left] at (11.62,-5.47) {VLA};
\node[num,anchor=east] at (14.06,-5.47) {18};
\draw[box,fill=white] (11.52,-6.12) rectangle (14.16,-7.01);
\node[lbl,anchor=west,align=left] at (11.62,-6.57) {teleop-system};
\node[num,anchor=east] at (14.06,-6.57) {14};
\draw[box,fill=white] (11.52,-7.17) rectangle (14.16,-8.05);
\node[lbl,anchor=west,align=left] at (11.62,-7.61) {diffusion};
\node[num,anchor=east] at (14.06,-7.61) {14};
\node[font=\footnotesize\bfseries,anchor=south] at (16.68,0.48) {evaluation};
\node[font=\scriptsize\color{black!60},anchor=south] at (16.68,0.10) {Sec.~\ref{sec:evaluation}};
\draw[draw=black!45,line width=0.4pt] (15.36,0.00) -- (18.00,0.00);
\draw[box,fill=white] (15.36,-0.00) rectangle (18.00,-2.95);
\node[lbl,anchor=west,align=left] at (15.46,-1.48) {real robot};
\node[num,anchor=east] at (17.90,-1.48) {89};
\draw[box,fill=white] (15.36,-3.11) rectangle (18.00,-4.24);
\node[lbl,anchor=west,align=left] at (15.46,-3.68) {simulation only};
\node[num,anchor=east] at (17.90,-3.68) {23};
\draw[box,fill=accent!8,draw=accent!60] (15.36,-4.40) rectangle (18.00,-5.21);
\node[lbl,anchor=west,align=left] at (15.46,-4.80) {penetration\\addressed};
\node[num,anchor=east] at (17.90,-4.80) {11};
\draw[box,fill=accent!8,draw=accent!60] (15.36,-5.37) rectangle (18.00,-8.21);
\node[lbl,anchor=west,align=left] at (15.46,-6.79) {penetration\\silent};
\node[num,anchor=east] at (17.90,-6.79) {85};
\end{tikzpicture}
}{%
    \framebox[0.7\textwidth]{\rule{0pt}{2.2cm}\textit{figs/fig\_field.tex pending}}}
  \caption{The field on one page. Data sources, embodiment, simulator, training paradigm and
  evaluation, with nodes sized by how many corpus papers sit in them and edges drawn only for
  routes the corpus actually contains. The routes converge: almost everything ends in a policy
  trained in a GPU simulator and distilled to vision.}
  \label{fig:field}
\end{figure*}
